\section{Cronograma e Próximos Passos}
\label{sec:cronograma}

\subsection{Cronograma de execução}
\label{sec:cron_exec}

O cronograma de execução da pesquisa abrange aproximadamente 28 a 32 semanas após a aprovação desta qualificação:

\begin{table}[h]
\centering
\small
\begin{tabular}{p{5cm}p{3cm}p{6cm}}
\toprule
\textbf{Fase} & \textbf{Target} & \textbf{Entrega} \\
\midrule
Aprovação da qualificação & Julho 2026 & Ajustes na proposta se necessário \\
Setup metodológico & Agosto 2026 & Documentos com especificações detalhadas \\
Cap. 1 --- Auditoria MST $\times$ race & Setembro 2026 & Resultado de H3 + protocolo OE2 \\
Cap. 2 --- Pipeline condicionado + ablation FiLM vs CLIP & Dezembro 2026 & Resultados H1, H2, H4 + C3, C7 \\
Cap. 3 --- Fair transferência para FR & Fevereiro 2027 & Resultados H5, H6 + C5 \\
Cap. 4 --- Síntese decompositiva & Março 2027 & Decomposição C6 \\
Escrita final & Em paralelo (Set 2026 -- Mar 2027) & Texto final da dissertação \\
Defesa & Abril 2027 & --- \\
\bottomrule
\end{tabular}
\caption{Cronograma de execução da pesquisa.}
\end{table}

\subsection{Riscos identificados e mitigações}
\label{sec:riscos}

\begin{description}[font=\bfseries\sffamily]
\item[Risco 1.] Alguma hipótese pode ser refutada. \\
\textbf{Mitigação:} Plano B documentado por hipótese; observação negativa é resultado científico válido.

\item[Risco 2.] Recrutamento de anotadores diversos para validação manual MST. \\
\textbf{Mitigação:} Prolific Academic com filtros regionais; custo estimado \$1.400 para 700 imagens $\times$ 3 anotadores.

\item[Risco 3.] Refutação de H5 por Pangelinan 2023 (pixel info $>$ skin tone em FR). \\
\textbf{Mitigação:} H5 reformulada com controle explícito de pixel info; H6 nova decompõe variância. Refutação parcial torna-se contribuição quantitativa.

\item[Risco 4.] Compute. \\
\textbf{Mitigação:} ConvNeXt-T é leve (28M parâmetros), aproximadamente 200--400 GPU-horas para o Capítulo 2 completo.
\end{description}

\subsection{Próximos passos imediatos pós-qualificação}
\label{sec:proximos}

\begin{enumerate}
\item Confirmar disponibilidade pública dos pesos do SkinToneNet \autocite{pereira2026skintonenet}; plano B: re-treinar ViT-Small em STW se necessário.
\item Recrutar anotadores via Prolific para validação manual MST em subset do FairFace.
\item Iniciar treinamento do Capítulo 1 (auditoria MST $\times$ race).
\item Implementar pipeline FiLM-conditioning sobre ConvNeXt-T \autocite{liu2022convnext}.
\item Implementar ablation arquitetural CLIP-conditioning \autocite{luo2024fairclip,radford2021clip}.
\end{enumerate}
